\chapter{シェルとファイル操作}
\label{chap:shell}

\section{ターミナル・シェル・WSL の役割}
\label{SEC:shell}
本書では、ファイル操作やスクリプト実行をターミナル\index{たーみなる@ターミナル|textbf}から行う。初めに、ターミナルとシェルの役割を区別する。

\begin{itemize}
\item ターミナル: 文字で命令を入力するためのアプリケーション
\item シェル\index{しぇる@シェル|textbf}: ターミナルの中で動き、入力されたコマンドを解釈して実行するプログラム
\end{itemize}

ターミナルを開くと、シェルはコマンドを入力できる位置に短い文字列を表示する。この文字列をプロンプト\index{ぷろんぷと@プロンプト|textbf}という。表示内容は環境によって異なるが、末尾には \code{\$} が使われることが多い。例えば、実際の画面には次のように表示される。

\begin{lstlisting}[numbers=none, xleftmargin=2\zw]
student@host ~ $
\end{lstlisting}

この表示は、シェルが入力を待っていることを意味する。本書の実行例では、実際のプロンプト全体を省略し、読者が入力する行の先頭に \code{\$} だけを置く。\code{\$} は入力行を見分けるための印であり、読者が入力する文字には含まれない。

ログインシェルとして設定されているシェルの名前は、次のコマンドで確認できる。

\begin{lstlisting}[numbers=none, xleftmargin=2\zw]
$ echo $SHELL
\end{lstlisting}

この実行例で読者が入力するのは、\code{echo \$SHELL} の部分だけである。行頭の \code{\$} は入力しない。

実行結果の例は次のようになる。

\begin{lstlisting}[numbers=none, xleftmargin=2\zw]
/bin/zsh
\end{lstlisting}

本書では、macOS の例に \code{zsh}、Linux と WSL 上の Ubuntu の例に \code{bash} を用いる。\code{zsh} と \code{bash} はどちらも本章で扱う基本コマンドを実行できるが、起動時に読む設定ファイルや一部のシェル構文は異なる。上の \code{/bin/zsh} は、ログインシェルとして \code{zsh} が設定されていることを表す。

ただし、\code{SHELL} は現在実行中のシェルを必ず示すわけではない。別のシェルを起動した場合などは値が一致しない。現在コマンドを解釈しているシェルは、次のコマンドで確認できる。

\begin{lstlisting}[numbers=none, xleftmargin=2\zw]
$ ps -p $$ -o comm=
zsh
\end{lstlisting}

設定ファイルは、シェル名だけでなく OS と起動方法によっても異なる。zsh の対話シェルでは通常 \code{\textasciitilde/.zshrc} を使う。bash では Linux の対話シェルなら \code{\textasciitilde/.bashrc}、macOS のログインシェルなら \code{\textasciitilde/.bash\_profile} が使われることが多い。

\subsection{Windows における WSL の利用}

本書のコマンド例は macOS や Linux のシェルを前提にしている。Windows では、WSL\index{wsl@WSL|textbf} (Windows Subsystem for Linux) 上の Ubuntu などを利用すると、同じコマンド例を実行しやすい。WSL は、管理者権限で開いた PowerShell から次のコマンドで導入できる。PowerShell のプロンプトは、本書の UNIX シェル用プロンプトと区別して示す。

\begin{lstlisting}[numbers=none, xleftmargin=2\zw]
PS C:\> wsl --install
\end{lstlisting}

WSL 上では、本書の Linux 用コマンドとディレクトリ構造を Windows 上で利用できる。また、後半で扱う並列処理や周辺ツールには Linux 系環境を前提とするものがあるため、Windows のコマンドプロンプトや PowerShell とは実行環境を分けて説明する。

\section{コマンド・オプション・引数とヘルプ}

\subsection{コマンド、オプション、引数}

シェルで打つ 1 行は、多くの場合「コマンド本体」「オプション」「引数」から成る。たとえば次のような形である。

\begin{lstlisting}[numbers=none, xleftmargin=2\zw]
$ ls -l results
$ cp data/raw.txt backup/raw.txt
\end{lstlisting}

1 行目では、\code{ls} がコマンド本体、\code{-l} がオプション、\code{results} が引数である。2 行目では、\code{cp} がコマンド本体で、\code{data/raw.txt} と \code{backup/raw.txt} が引数である。一般に、オプションは動作を変え、引数は対象のファイル名やディレクトリ名を与える。

多くのコマンドには、\code{-} と 1 文字からなる短いオプション、または \code{--} と英単語からなる長いオプションがある。ただし、オプション名と動作はコマンドごとに定義される。

\subsection{\code{--help} と \code{man} による使用法の確認}

\code{--help} は単独のコマンドではなく、対象のコマンド名に続けて指定するオプションである。例えば \code{python3 --help} は、\code{python3} が受け付けるオプションと引数の書式を表示する。\code{python3} では \code{-h} も同じ説明を表示する。

\code{man} は、コンピュータに収録されている説明書を表示するコマンドである。\code{man} の後ろに調べたいコマンド名を置く。例えば \code{man rsync} を実行すると、\code{rsync} の説明書が開く。実際の入力例は次のとおりである。

\begin{lstlisting}[numbers=none, xleftmargin=2\zw]
$ python3 --help
$ python3 -h
$ man rsync
\end{lstlisting}

\code{python3 -h} の出力は長いが、先頭はたとえば次のようになる。

\begin{lstlisting}[numbers=none, xleftmargin=2\zw]
usage: python3 [option] ... [-c cmd | -m mod | file | -] [arg] ...
Options and arguments (and corresponding environment variables):
-b     : issue warnings about str(bytes_instance), str(bytearray_instance)
         and comparing bytes/bytearray with str. (-bb: issue errors)
-B     : don't write .pyc files on import; also PYTHONDONTWRITEBYTECODE=x
-c cmd : program passed in as string (terminates option list)
\end{lstlisting}

ここで最初の \code{usage: ...} は「このコマンドをどういう形で呼び出せるか」を要約した行である。\code{[-c cmd | -m mod | file | -]} のように角括弧で囲まれた部分は省略可能、\code{|} で区切られた部分はどれか 1 つを選ぶ、という意味で読む。その次の \code{Options and arguments ...} 以下では、各オプションの意味が 1 行ずつ並ぶ。\code{-b} のように短いオプション名、コロンの後ろにその効果、必要なら対応する環境変数が示される。\code{-c cmd} のように引数を伴うものは、オプションの直後に必要な値の名前も書かれる。

\code{man rsync} を開くと、画面の先頭付近はたとえば次のようになる。

\begin{lstlisting}[numbers=none, xleftmargin=2\zw]
OPENRSYNC(1)                                    General Commands Manual                                   OPENRSYNC(1)

NAME
     openrsync - synchronise local and remote files

SYNOPSIS
     openrsync [-0468BCDEFHILOPRSTWVabcdefghiklmnopqrtuvxyz] [-e program] [-f filter] ...
               ... source ... directory
\end{lstlisting}

\code{NAME} はコマンド名と機能の要約、\code{SYNOPSIS} は呼び出し書式を示す。角括弧 \code{[...]} 内は省略可能である。この例の末尾にある \code{source ... directory} は、一つ以上のコピー元に続けてコピー先を書く構造を表す。必要なオプションの綴りと意味は、後続の説明または \code{man} で確認する。

\code{--help} や \code{-h} は、そのコマンドが受け付けるオプションや引数の書式を短く表示することが多い。ただし、どのコマンドでも両方を使えるとは限らない。例えば macOS の \code{ls} は \code{--help} を受け付けない。その場合は、\code{man ls}、\code{man ssh}、\code{man rsync} のように、\code{man} の後ろへ調べたいコマンド名を置く。
\code{man} を開いたあとは、上下キーや \code{Space} で移動し、\code{/pattern} で検索できる。\code{q} を押せば終了する。

検索エンジンや AI を用いてコマンドの使い方を調べることもできるが、同じ名前のコマンドでもバージョンや OS によって使い方が異なる場合がある。インターネット上の説明と手元の動作が食い違う場合は、手元のコマンドに \code{--help} を付けた出力、または \code{man} で開いた説明書を基準とする。

\section{現在位置・一覧・パス・ディレクトリ操作}

\subsection{現在位置の確認と項目の一覧: \code{pwd} と \code{ls}}

シェルでは、現在作業しているディレクトリがパス解釈の基準となる。このディレクトリをカレントディレクトリという。現在位置は \code{pwd}、その場所にある項目は \code{ls} で確認できる。

\begin{lstlisting}[numbers=none, xleftmargin=2\zw]
$ pwd
$ ls
$ ls -l
$ ls -lh
$ ls -a
$ ls -lah
\end{lstlisting}

\code{pwd}\index{pwd@\code{pwd}|textbf} は、現在のカレントディレクトリの絶対パスを表示する。\code{ls}\index{ls@\code{ls}|textbf} は、そのディレクトリにあるファイルやディレクトリの一覧を表示する。たとえば \code{pwd} の結果が次のようであれば、

\begin{lstlisting}[numbers=none, xleftmargin=2\zw]
/Users/student/analysis
\end{lstlisting}

現在の作業場所は \code{/Users/student/analysis} である。この状態で \code{ls} を実行すると、その中にある項目が表示される。

\begin{lstlisting}[numbers=none, xleftmargin=2\zw]
$ ls
data	scripts
\end{lstlisting}

この結果から、\code{/Users/student/analysis} の中には \code{data} と \code{scripts} という 2 つのディレクトリがあることが分かる。

\code{ls} の表示内容はオプションで変更できる。\code{-l} は、各項目の種別、アクセス権、リンク数、所有者、グループ、サイズ、最終更新日時、名前を表示する。

\begin{lstlisting}[numbers=none, xleftmargin=2\zw]
$ ls -l
total 0
drwxr-xr-x  3 student  staff  96  4月 12 23:19 data
drwxr-xr-x  3 student  staff  96  4月 12 23:19 scripts
\end{lstlisting}

\code{ls -l} の先頭にある \code{total} は、表示対象へ割り当てられたファイルシステムブロック数の合計である。1 ブロックを表示上何バイトと数えるかは実装や環境変数で異なるため、ファイルごとのバイト数には各行のサイズ欄、ディスク使用量には \code{du} を用いる。

\code{ls -l} の出力では、その下の各行が1つのファイルまたはディレクトリに対応しており、左から順に次の情報が表示される。

\begin{itemize}
  \item \code{drwxr-xr-x}：ファイル種別とアクセス権
  \item \code{3}：リンク数
  \item \code{student}：所有者
  \item \code{staff}：グループ
  \item \code{96}：サイズ
  \item \code{4月 12 23:19}：最終更新日時
  \item \code{data}：名前
\end{itemize}

先頭の1文字は種別を表し、\code{d} ならディレクトリ、\code{-} なら通常のファイルである。続く \code{rwxr-xr-x} はアクセス権を表している。アクセス権の詳しい意味は第~\ref{subsec:mod}~節で説明する。

リンク数は、同じ inode を参照するハードリンクの数を表す。通常ファイルを新規作成した直後は一般に 1 であり、\code{ln} でハードリンクを追加すると増える。ディレクトリのリンク数はファイルシステムとサブディレクトリ構造に依存し、ディレクトリ内の項目数そのものではない。

\code{ls -l} のサイズ欄はバイト単位である。\code{-h} を併用すると、同じ値が \code{K}、\code{M}、\code{G} などの接頭辞付きで表示される。表示の丸めを避けて正確なバイト数を記録する場合は、\code{-h} を付けない出力を用いる。

\begin{lstlisting}[numbers=none, xleftmargin=2\zw]
$ ls -lh
total 0
drwxr-xr-x  3 student  staff    96B  4月 12 23:19 data
drwxr-xr-x  3 student  staff    96B  4月 12 23:19 scripts
\end{lstlisting}

この例の \code{96B} はディレクトリ項目そのものの大きさであり、配下のファイルサイズの合計ではない。ディレクトリ以下のディスク使用量を調べる場合は、\code{du} を用いる。

通常のファイルに対して \code{ls -lh} を使うと、バイト数は \code{12K} や \code{3.4M} のように SI 接頭語に似た接尾辞付きで表示される。正確なバイト数が必要な場合は、OS に応じて \code{stat} または \code{wc -c} を用いる。

一方、\code{-a} オプションを付けると、ドットで始まる隠し項目も含めて表示できる。

\begin{lstlisting}[numbers=none, xleftmargin=2\zw]
$ ls -a
.	..	.config	data	scripts
\end{lstlisting}

名前が \code{.}（ドット）で始まるファイルやディレクトリは隠し項目として扱われ、通常の \code{ls} では表示されない。たとえば \code{.config}、\code{.gitignore}、\code{.bashrc} などがこれにあたる。設定ファイルなど、通常の一覧に常時表示する必要のない項目にこの命名が用いられる。

上の \code{ls -a} に現れる \code{.} と \code{..} にも意味がある。\code{.} は現在のディレクトリ自身を表し、\code{..} は1つ上の親ディレクトリを表す。これらは後でディレクトリ移動をするときにも重要になる。

\code{-lah} は、詳細表示の \code{-l}、隠し項目を含める \code{-a}、サイズに接尾辞を付ける \code{-h} を組み合わせた指定である。

\begin{lstlisting}[numbers=none, xleftmargin=2\zw]
$ ls -lah
total 96
drwxr-xr-x  5 student  staff   160B  4月 12 23:20 .
drwxr-xr-x  5 student  staff   160B  4月 12 23:18 ..
-rw-r--r--  1 student  staff    47K  4月 12 23:22 .config
drwxr-xr-x  3 student  staff    96B  4月 12 23:19 data
drwxr-xr-x  3 student  staff    96B  4月 12 23:19 scripts
\end{lstlisting}

この例では \code{.config} という通常のファイルも表示されており、\code{ls -l} の場合とは \code{total} の値が異なる。\code{total} は各行のサイズ欄の合計ではなく、表示対象へ割り当てられたファイルシステムブロック数に基づく。

\code{pwd} は現在位置、\code{ls} は項目名、\code{ls -l} は権限や更新時刻を含む詳細、\code{ls -a} は隠し項目を含む一覧を表示する。\code{-h} を \code{-l} と組み合わせると、ファイル容量が単位付きで表示される。必要な情報に応じて、これらのオプションを選択する。

\subsection{ディレクトリの移動: \code{cd}}

\code{cd}\index{cd@\code{cd}|textbf} はカレントディレクトリを変更するコマンドである。\code{ls}、\code{cat}、\code{cp} などへ相対パスを渡すと、そのパスはカレントディレクトリを基準に解釈される。次のディレクトリ構成を例に、移動前後でパスの指す対象を確認する。

\begin{lstlisting}[numbers=none, xleftmargin=2\zw]
analysis/
|-- data/
|   `-- events.txt
`-- scripts/
    `-- plot.py
\end{lstlisting}

いま \code{/Users/student/analysis} にいるとする。まず \code{pwd} と \code{ls} を実行すると、現在位置とその場にある項目は次のように分かる。

\begin{lstlisting}[numbers=none, xleftmargin=2\zw]
$ pwd
/Users/student/analysis
$ ls
data	scripts
\end{lstlisting}

この状態で \code{cd data} を実行すると、\code{analysis} の下にある \code{data} ディレクトリへ移動する。

\begin{lstlisting}[numbers=none, xleftmargin=2\zw]
$ cd data
$ pwd
/Users/student/analysis/data
$ ls
events.txt
\end{lstlisting}

\code{data} は相対パスなので、カレントディレクトリ \code{/Users/student/analysis} を基準とする \code{/Users/student/analysis/data} を指す。

次に \code{cd ..} を実行すると、1 つ上の親ディレクトリへ戻る。

\begin{lstlisting}[numbers=none, xleftmargin=2\zw]
$ cd ..
$ pwd
/Users/student/analysis
\end{lstlisting}

ここで使っている \code{..} は、\code{ls -a} でも見た「親ディレクトリ」を表す特別な名前である。

カレントディレクトリに依存しない指定で直接移動することもできる。\code{\textasciitilde} は、シェルがホームディレクトリのパスへ展開する記号である。次の指定では、現在位置に関係なく \code{scripts} へ移動できる。

\begin{lstlisting}[numbers=none, xleftmargin=2\zw]
$ cd ~/analysis/scripts
$ pwd
/Users/student/analysis/scripts
$ ls
plot.py
\end{lstlisting}

また、\code{cd -} は 1 つ前にいた場所へ戻る。たとえば上の状態で実行すると、直前にいた \code{/Users/student/analysis} へ戻る。

\begin{lstlisting}[numbers=none, xleftmargin=2\zw]
$ cd -
/Users/student/analysis
\end{lstlisting}

引数なしの \code{cd} は、現在位置にかかわらずホームディレクトリへ移動する。

\begin{lstlisting}[numbers=none, xleftmargin=2\zw]
$ cd
$ pwd
/Users/student
\end{lstlisting}

\code{cd} には、相対パスによる移動、絶対パスによる移動、\code{..} による親ディレクトリへの移動、\code{-} による直前のディレクトリへの移動がある。移動後に \code{pwd} と \code{ls} を実行すれば、現在位置とその内容を確認できる。

\subsection{絶対パスと相対パス}

ファイルやディレクトリの場所は、パス\index{ぱす@パス|textbf}で表す。パスは階層を \code{/} でつないだ表記であり、最初の書き方によって意味が変わる。大きく分けると絶対パスと相対パスがある。

\begin{table}[hbtp]
\centering
\caption{よく出てくるパス表記}
\label{tab:path-notation}
\begin{tabular}{ll}
\toprule
表記 & 意味 \\
\midrule
\code{/abs/data.txt} & 絶対パス。\code{/} から始まり、どこにいても同じ場所を指す \\
\code{\textasciitilde/analysis} & ホームディレクトリからのパス。実行前にシェルが展開する \\
\code{data.txt} & 相対パス。今いるディレクトリを基準に解釈される \\
\code{./x.txt} & 相対パス。\code{.} は現在のディレクトリ \\
\code{../x.txt} & 相対パス。\code{..} は 1 つ上のディレクトリ \\
\bottomrule
\end{tabular}
\end{table}

表~\ref{tab:path-notation} は、シェルを読み始めた直後に最低限区別したい表記をまとめたものである。絶対パスは \code{/} から始まる。\code{\textasciitilde/analysis} はシェルによる展開後にホームディレクトリ下の絶対パスとなり、カレントディレクトリには依存しない。

次のディレクトリ構成を考える。

\begin{lstlisting}[numbers=none, xleftmargin=2\zw]
analysis/
|-- data/
|   `-- events.txt
`-- scripts/
    `-- plot.py
\end{lstlisting}

現在位置が \code{/Users/student/analysis} なら、次の 2 つのコマンドは同じファイルを指している。

\begin{lstlisting}[numbers=none, xleftmargin=2\zw]
$ pwd
/Users/student/analysis
$ cat data/events.txt
$ cat /Users/student/analysis/data/events.txt
\end{lstlisting}

1 つ目の \code{data/events.txt} は、カレントディレクトリ \code{/Users/student/analysis} を基準とする相対パスである。2 つ目の \code{/Users/student/analysis/data/events.txt} は絶対パスであり、カレントディレクトリに依存しない。この状態では、両者とも同じファイルを指す。

一方で、現在位置が \code{/Users/student/analysis/scripts} に変わると、同じ相対パスでも意味が変わる。

\begin{lstlisting}[numbers=none, xleftmargin=2\zw]
$ pwd
/Users/student/analysis/scripts
$ cat data/events.txt
cat: data/events.txt: No such file or directory
$ cat ../data/events.txt
\end{lstlisting}

この場所で \code{data/events.txt} と指定すると、シェルは \code{scripts/data/events.txt} を探す。そのファイルは存在しないため、\code{data/events.txt: No such file or directory} というエラー文が表示される。
相対パスで \code{events.txt} を指定するには、1 つ上へ戻る \code{..} を使って \code{../data/events.txt} と書く必要がある。

同様に、\code{./plot.py} は現在のディレクトリにある \code{plot.py}、\code{../data} は一つ上のディレクトリにある \code{data} を表す。相対パスが指す対象はカレントディレクトリによって変わるため、コマンドの実行前に \code{pwd} と対象パスを確認する。

\subsection{ディレクトリの作成: \code{mkdir}}

\code{mkdir}\index{mkdir@\code{mkdir}|textbf} は新しいディレクトリを作るコマンドである。解析では、生データ、途中結果、図、ログなどを別々の場所へ整理したいので、作業の最初に出力先のディレクトリを作ることが多い。

たとえば、最初に次のような構成になっているとする。

\begin{lstlisting}[numbers=none, xleftmargin=2\zw]
analysis/
`-- data/
    `-- events.txt
\end{lstlisting}

\code{output} というディレクトリを 1 つ作るだけなら、次のように実行すればよい。

\begin{lstlisting}[numbers=none, xleftmargin=2\zw]
$ mkdir output
$ ls
data	output
\end{lstlisting}

1 行目では、現在のカレントディレクトリの直下に \code{output} が新しくできる。2 行目の \code{ls} で、それが確認できる。

一方、\code{output/2026/04} のような深い階層をまとめて作る場合もある。途中の \code{output/2026} がまだ無ければ、\code{mkdir output/2026/04} だけではエラーになる。\code{-p} オプションを付けると、存在しない親ディレクトリも含めて作成される。

\begin{lstlisting}[numbers=none, xleftmargin=2\zw]
$ mkdir output/2026/04
mkdir: output/2026: No such file or directory
$ mkdir -p output/2026/04
\end{lstlisting}

実行後は、次のようなディレクトリ構成になる。

\begin{lstlisting}[numbers=none, xleftmargin=2\zw]
analysis/
|-- data/
|   `-- events.txt
`-- output/
    `-- 2026/
        `-- 04/
\end{lstlisting}

観測日ごとの出力先や \code{figures/summary/} のような階層を一度に作る場合は、\code{mkdir -p} を用いる。\code{-p} は必要な親ディレクトリも作成し、指定したディレクトリがすでに存在する場合にはエラーにしない。

\section{コピー、移動、削除}

\subsection{ファイル・ディレクトリのコピー: \code{cp}}

\code{cp}\index{cp@\code{cp}|textbf} は元のファイルやディレクトリを残したまま複製を作るコマンドである。最後の引数が既存のディレクトリならその中へ同じ名前でコピーされ、ファイル名ならその名前で複製が作られる。以下のディレクトリ構造を例とする。

\begin{lstlisting}[numbers=none, xleftmargin=2\zw]
work/
|-- backup/
`-- raw/
    `-- events.txt
\end{lstlisting}

ここで \code{raw/events.txt} を \code{backup} へコピーするには、たとえば次のように実行できる。

\begin{lstlisting}[numbers=none, xleftmargin=2\zw]
$ cp raw/events.txt backup/
$ cp raw/events.txt backup/events_20260412.txt
\end{lstlisting}

1 行目では、\code{backup/} が既存のディレクトリなので、その中に同じ名前の \code{events.txt} が作られる。2 行目では、コピー先の名前まで明示しているので、同じ内容を \code{events\_20260412.txt} という別名で保存できる。

\begin{lstlisting}[numbers=none, xleftmargin=2\zw]
work/
|-- backup/
|   |-- events.txt
|   `-- events_20260412.txt
`-- raw/
    `-- events.txt
\end{lstlisting}

\code{cp} の実行後も \code{raw/events.txt} は残り、コピー先に独立したファイルが作成される。元データを保持したまま加工用ファイルを作る場合は、入力と出力のパスを分ける。

ディレクトリごと複製したい場合は \code{-r} を付ける。

\begin{lstlisting}[numbers=none, xleftmargin=2\zw]
$ cp -r raw raw_backup
\end{lstlisting}

\begin{lstlisting}[numbers=none, xleftmargin=2\zw]
work/
|-- backup/
|   |-- events.txt
|   `-- events_20260412.txt
|-- raw/
|   `-- events.txt
`-- raw_backup/
    `-- events.txt
\end{lstlisting}

\code{-r} は、指定したディレクトリ以下を再帰的にコピーする。コピー先に同名の項目が存在する可能性がある場合、\code{-i} を指定すると上書き前に確認が表示される。自動処理では対話入力に依存させず、出力先の存在を事前に検査する。

\subsection{移動と名前の変更: \code{mv}}

\code{mv}\index{mv@\code{mv}|textbf} は「別の場所へ移す」と「名前を変える」の両方に使う。考え方は \code{cp} と似ており、最後の引数が既存ディレクトリならその中へ移動し、そうでなければその名前へ変更する。ただし、\code{cp} と違って元の場所には残らない。

まず、同じディレクトリの中で名前だけを変える例を見る。

\begin{lstlisting}[numbers=none, xleftmargin=2\zw]
work/
`-- result.txt
\end{lstlisting}

\begin{lstlisting}[numbers=none, xleftmargin=2\zw]
$ mv result.txt result_old.txt
\end{lstlisting}

\begin{lstlisting}[numbers=none, xleftmargin=2\zw]
work/
`-- result_old.txt
\end{lstlisting}

この場合は、同じ場所にあるファイルの名前だけが \code{result\_old.txt} へ変わる。

次に、別のディレクトリへ移動しながら名前も変える例を考える。

\begin{lstlisting}[numbers=none, xleftmargin=2\zw]
work/
|-- output/
`-- result_old.txt
\end{lstlisting}

\begin{lstlisting}[numbers=none, xleftmargin=2\zw]
$ mv result_old.txt output/final.txt
\end{lstlisting}

\begin{lstlisting}[numbers=none, xleftmargin=2\zw]
work/
`-- output/
    `-- final.txt
\end{lstlisting}

\code{mv result\_old.txt output/final.txt} では、\code{result\_old.txt} が \code{output} の中へ移り、名前も \code{final.txt} に変わる。もし \code{mv result\_old.txt output/} と書けば、名前はそのままで場所だけが変わる。

このように、\code{mv} は「移動」と「改名」を 1 つのコマンドで行えるのが特徴である。整理のためにファイルを別ディレクトリへ移したり、解析段階に応じて \code{tmp.txt} を \code{final.txt} に改名したりするときに使う。\code{mv -i} を使うと、すでにファイルが存在した場合に上書きするかどうかを確認しながら進められる。

\subsection{ファイル・ディレクトリの削除: \code{rm}}

\code{rm}\index{rm@\code{rm}|textbf} はファイルやディレクトリを削除するコマンドである。通常の \code{rm} は対象をゴミ箱へ移さず、その場で削除する。そのため、実行前に \code{pwd} や \code{ls} で場所と対象を確認する。

たとえば、次のような状態を考える。

\begin{lstlisting}[numbers=none, xleftmargin=2\zw]
work/
|-- old_output/
|   `-- plot.png
`-- result_old.txt
\end{lstlisting}

ここで使う基本形は次の通りである。

\begin{lstlisting}[numbers=none, xleftmargin=2\zw]
$ rm -i result_old.txt
$ rm -ri old_output
\end{lstlisting}

\code{rm -i result\_old.txt} は通常のファイルを削除する前に確認を表示する。\code{rm -ri old\_output} は、ディレクトリの内部を再帰的にたどり、各対象の削除前に確認する。確認が不要で対象を確定できている場合は、それぞれ \code{-i} を外した \code{rm result\_old.txt} または \code{rm -r old\_output} を用いる。これらは別々の初期状態に対する代替例であり、同じ対象に順番に実行するものではない。

\code{rm -r} は、指定したディレクトリ以下を再帰的に削除する。対象パスの誤指定による削除範囲の拡大を防ぐため、実行前に \code{ls old\_output} や \code{ls result\_old.txt} で対象を確認する。内容を個別に確認する場合は \code{rm -ri} を用いる。

\subsection{ワイルドカード}

ワイルドカードは、パターンに一致する複数のパスを指定する記法である。\code{*} は、ディレクトリ区切りを除く 0 文字以上の文字列に一致する。\code{cat} や \code{rm} の実行前にシェルがパターンを実際のパスへ展開し、その結果を各コマンドの引数として渡す。

たとえば、次のような構成を考える。

\begin{lstlisting}[numbers=none, xleftmargin=2\zw]
data/
|-- run01.csv
|-- run02.csv
|-- run03.txt
|-- test_20260411.csv
|-- test_20260412.csv
|-- test_20260413.csv
`-- subdir/
    `-- run04.csv
\end{lstlisting}

このとき、次のような書き方ができる。

\begin{lstlisting}[numbers=none, xleftmargin=2\zw]
$ ls data/*.csv
$ cp data/*.csv backup/
$ rm data/test_*.csv
\end{lstlisting}

\code{data/*.csv} は、\code{data} ディレクトリ直下にある \code{.csv} ファイルすべてを意味する。したがって、この例では \code{run01.csv} と \code{run02.csv} は対象になるが、\code{subdir/run04.csv} は含まれない。\code{*} は \code{/} をまたいで階層全体を探すわけではないからである。

\code{cp data/*.csv backup/} とすると、条件に合う複数の CSV をまとめてコピーできる。\code{rm data/test\_*.csv} なら、\code{data/test\_001.csv} や \code{data/test\_20260413.csv} のように同じ接頭辞を持つファイルを一括削除できる。

ただし、ワイルドカード\index{わいるどかーど@ワイルドカード|textbf}は想定より多くのファイルへ一致することがある。特に削除系のコマンドでは、先に \code{ls data/test\_*.csv} で実際の対象を表示してから \code{rm data/test\_*.csv} を実行する。また、一致する項目がない場合の振る舞いはシェル設定によって異なるため、削除前の一覧表示はその確認にもなる。


\section{ファイルの内容確認とデータ抽出}

データを処理する前には、ファイル形式と内容の一部を確認する。表~\ref{tab:file-inspection-commands} に、内容表示、行抽出、列処理に用いるコマンドと主要なオプションを示す。

\begin{table}[hbtp]
\centering
\caption{内容表示・行抽出・列処理のコマンドと主要オプション}
\label{tab:file-inspection-commands}
\begin{tabular}{p{0.14\linewidth}p{0.28\linewidth}p{0.44\linewidth}}
\toprule
コマンド & 主なオプションや書式 & 主な用途 \\
\midrule
\code{cat} & \code{-n} & 短いファイルをそのまま読む、複数ファイルを連結する \\
\code{less} & \code{-N}, \code{/pattern}, \code{q} & 長いファイルをスクロールしながら読む \\
\code{head} & \code{-n N} & 先頭の数行だけを見る \\
\code{tail} & \code{-n N}, \code{-f} & 末尾の数行を見る、ログの追記を監視する \\
\code{wc} & \code{-l}, \code{-w}, \code{-c} & 行数、単語数、バイト数を数える \\
\code{grep} & \code{-i}, \code{-v}, \code{-n}, \code{-A N}, \code{-B N} & 条件に合う行だけを抽出する \\
\code{awk} & \code{-F}, \code{if (...)}, \code{sum += ...} & 列単位の抽出、整形、集計を行う \\
\bottomrule
\end{tabular}
\end{table}

表~\ref{tab:file-inspection-commands} は、ファイルの確認と抽出に用いるコマンドを役割別に示したものである。「内容の表示」「条件に合う行の抽出」「列単位の整形」を区別し、必要な操作に対応するコマンドを選ぶ。以下では、各コマンドの入力、出力、主要なオプションを順に説明する。

\subsection{\code{cat}}

\code{cat}\index{cat@\code{cat}|textbf} は、指定したファイルの内容を順に標準出力へ送る。複数のファイルを渡すと、引数の順に内容を連結する。

\begin{lstlisting}[numbers=none, xleftmargin=2\zw]
$ cat data/events.txt
$ cat header.txt body.txt
$ cat log/*.log
$ cat -n config.txt
\end{lstlisting}

3 行目のように、ワイルドカードを使えば複数のログファイルを連結して表示できる。ただし、長いログや観測データでは出力が画面の範囲を超える。短いファイルの全体表示には \code{cat}、長いファイルの閲覧や一部表示には \code{less}、\code{head}、\code{tail} を用いる。

\begin{description}
\item[\code{-n}] 各行に行番号を付ける。設定ファイルやログの位置を行番号で参照できる。
\end{description}

\subsection{\code{less}}

\code{less}\index{less@\code{less}|textbf} は、ファイルを画面単位で表示し、開いたまま移動と検索を行うコマンドである。

\begin{lstlisting}[numbers=none, xleftmargin=2\zw]
$ less data/events.txt
$ less -N data/events.txt
\end{lstlisting}

\code{less} では、ファイルを開いたまま検索と移動を行える。主要な操作を次に示す。

\begin{description}
\item[\code{-N}] 行番号を表示する。後で「どの行に何があったか」を参照しやすくなる。
\item[\code{Space}] 1 画面ぶん先へ進む。
\item[\code{b}] 1 画面ぶん戻る。
\item[\code{/pattern}] 指定した文字列を前方検索する。たとえば \code{/ERROR} と入力すると、\code{ERROR} を含む次の行が表示位置となる。
\item[\code{n}] 直前の検索条件の次の一致へ進む。
\item[\code{q}] 表示を終了してターミナルへ戻る。
\end{description}

\subsection{\code{head}}

\code{head}\index{head@\code{head}|textbf} は、ファイルの先頭から指定した行数を表示する。ヘッダ、コメント行、区切り文字の確認に用いる。

\begin{lstlisting}[numbers=none, xleftmargin=2\zw]
$ head data/events.txt
$ head -n 20 data/events.txt
\end{lstlisting}

\begin{description}
\item[\code{-n N}] 先頭から \code{N} 行を表示する。指定しない場合は 10 行である。
\end{description}

\subsection{\code{tail}}

\code{tail}\index{tail@\code{tail}|textbf} はファイル末尾の行を表示する。ログの直近部分の確認や、末尾へ追記される計算結果の監視に用いる。

\begin{lstlisting}[numbers=none, xleftmargin=2\zw]
$ tail log.txt
$ tail -n 30 log.txt
$ tail -f log.txt
\end{lstlisting}

\begin{description}
\item[\code{-n N}] 末尾から \code{N} 行を表示する。
\item[\code{-f}] ファイルの末尾を監視し、追記された内容を継続して表示する。実行中の処理がログへ追記する内容を確認できる。
\end{description}

\subsection{\code{wc}}

\code{wc}\index{wc@\code{wc}|textbf} は、ファイルの行数、単語数、バイト数を数える。1 行が 1 レコードに対応する形式では、\code{-l} の結果からヘッダ行とコメント行を除くことでレコード数を確認できる。

\begin{lstlisting}[numbers=none, xleftmargin=2\zw]
$ wc data/events.txt
$ wc -l data/events.txt
$ wc -w memo.txt
$ wc -c output.bin
\end{lstlisting}

\begin{description}
\item[\code{-l}] 改行の数を行数として表示する。末尾に改行がない最終行の扱いに注意する。
\item[\code{-w}] 空白で区切られた単語数を表示する。
\item[\code{-c}] 入力のバイト数を表示する。
\end{description}

\subsection{\code{grep}}

\code{grep}\index{grep@\code{grep}|textbf} は、指定したパターンに一致する行を出力する。大文字・小文字の無視、非一致行の出力、行番号や前後行の付加をオプションで指定できる。

\begin{lstlisting}[numbers=none, xleftmargin=2\zw]
$ grep "ERROR" log.txt
$ grep -n "ERROR" log.txt
$ grep -A 2 -B 1 "ERROR" log.txt
$ grep -i "warning" log.txt
$ grep -v "^#" data.txt
$ grep -rn "TODO" src/
\end{lstlisting}

\begin{description}
\item[\code{-i}] 大文字と小文字を区別せずに検索する。たとえば、\code{error} と \code{ERROR} の両方に一致する。
\item[\code{-v}] 一致しなかった行を出す。コメント行や除外対象の行を取り除ける。
\item[\code{-n}] 一致した行とともに、元ファイル内の行番号を表示する。
\item[\code{-A N}] 一致した行に続く \code{N} 行も表示する。エラー行と直後の補足メッセージを同時に確認できる。
\item[\code{-B N}] 一致した行より前の \code{N} 行も表示する。エラー行に先行するログを同時に確認できる。
\item[\code{-r}] 指定したディレクトリ以下を再帰的に検索する。\code{grep -rn "TODO" src/} は、\code{src/} 以下で \code{TODO} を含む行を、ファイル名と行番号付きで表示する。
\end{description}

\code{grep} では正規表現も使用できる。\verb|^| は行頭、\verb|$| は行末を表すため、\verb|grep "^#" data.txt| は \code{\#} で始まる行だけを抽出する。\code{grep -r "ERROR" log/} は、\code{log/} 以下を再帰的に検索する。\verb|[]|、\verb|.|、\verb|*| を含む正規表現の記法は、付録~\ref{app:regex} で説明する。

\subsection{\code{awk}}

\code{awk}\index{awk@\code{awk}|textbf} は、入力行を区切り文字で列へ分け、列条件による抽出、列の並べ替え、集計を行うコマンドである。既定の区切りは連続する空白であり、\code{-F} で別の区切り文字を指定できる。

\begin{lstlisting}[numbers=none, xleftmargin=2\zw]
$ awk '{print $1}' data.txt
$ awk '{print $1, $3}' data.txt
$ awk -F, '{print $2}' data.csv
$ awk '{ if ($3 > 10) print $1, $3 }' data.txt
$ awk '{ sum += $3 } END { print sum }' data.txt
\end{lstlisting}

既定では空白やタブが区切り文字として使われる。そのため、空白区切りの観測データでは \code{\$1} が 1 列目、\code{\$2} が 2 列目を表す。条件分岐を \code{if} で明示すれば、「3 列目が 10 より大きい行だけを表示する」という処理を記述できる。各行の値を加算し、最後に合計や平均を出力することもできる。

\begin{description}
\item[\code{-F}] 区切り文字を指定する。CSV のようにカンマ区切りなら \code{-F,} のように書く。
\item[\code{\{print \$1\}}] 1 列目を表示する最も基本的な形である。\code{\{print \$1, \$3\}} のようにすれば複数列を並べて出せる。
\item[\code{\{ if (\$3 > 10) print \$1, \$3 \}}] 3 列目が 10 より大きい行について、1 列目と 3 列目を表示する。
\item[\code{\{ sum += \$3 \} END \{ print sum \}}] 各行で 3 列目を足し上げ、入力をすべて処理した後に合計を表示する。
\end{description}

\subsection{パイプ \code{|}}

パイプ\index{ぱいぷ@パイプ|textbf}（\code{|}）は、左側のコマンドの標準出力を、右側のコマンドの標準入力へ接続する。中間ファイルを作らずに、抽出、並べ替え、行数の集計などを連結できる。

\begin{lstlisting}[numbers=none, xleftmargin=2\zw]
$ grep "ERROR" log.txt | head -n 10
$ grep "ERROR" log.txt | wc -l
$ grep -v "noise" events.txt | awk '{print $2}'
\end{lstlisting}

1 行目は \code{ERROR} を含む行を抽出し、その先頭 10 行だけを表示する。2 行目は同じ抽出結果を \code{wc -l} へ渡し、件数を数える。3 行目はノイズ行を除外したうえで、2 列目だけを取り出す。

パイプによって、行の抽出、件数の集計、列の選択を左から右の順に一つのコマンド列として記述できる。

\section{\code{curl}・\code{wget} による URL からのファイル取得}

解析では、公開データ、設定例、インストール用スクリプトを URL から取得することがある。以下では、受信内容を既定で標準出力へ送る \code{curl} と、URL の末尾から決めた名前でファイルへ保存する \code{wget} を扱う。macOS の標準環境では \code{curl} を利用できるが、\code{wget} の利用には Homebrew などによる追加が必要である。

\subsection{\code{curl}}

\code{curl}\index{curl@\code{curl}|textbf} は、URL から受信した内容を標準出力へ出すか、指定したファイルへ保存する。標準出力、URL 末尾の名前、指定した名前という 3 通りの出力先を次に示す。

\begin{lstlisting}[numbers=none, xleftmargin=2\zw]
$ curl https://example.com/config.txt | head -n 2
$ curl -O https://example.com/data/run001.csv
$ curl -o events.csv https://example.com/data/run001.csv
$ curl -L -O https://example.com/latest.csv
$ curl -fsSL -o pyenv-installer.sh https://pyenv.run
$ less pyenv-installer.sh
\end{lstlisting}

1 行目では、取得した内容をファイルへ保存せず、そのまま標準出力へ送っている。配布元の \code{config.txt} の先頭が次の 2 行であれば、同じ内容が表示される。

\begin{lstlisting}[numbers=none, xleftmargin=2\zw]
threshold=5.0
window_ns=250
\end{lstlisting}

2 行目と 3 行目はファイル保存である。2 行目の \code{-O} は URL の末尾にあるファイル名をそのまま使う。3 行目の \code{-o events.csv} は、保存名を自分で決めたいときに使う。4 行目の \code{-L} はリダイレクト先までたどるオプションである。5 行目は取得したスクリプトをファイルへ保存し、6 行目で実行前に内容を確認する例である。取得内容を直接シェルへ渡すと、確認していないコードが実行されるため、配布元と内容の確認が必要である。\code{-fsSL} の各文字は、\code{-f} が HTTP エラーを失敗として扱い、\code{-sS} が通常の進捗を抑制しつつエラーを表示し、\code{-L} がリダイレクトをたどることを表す。

\begin{description}
\item[\code{-O}] URL の末尾にある名前で保存する。保存名を URL と一致させる場合に用いる。
\item[\code{-o}] 保存名を引数で指定する。URL から保存名を決められない場合にも利用できる。
\item[\code{-L}] HTTP のリダイレクト応答で示された URL をたどる。リダイレクトを返す URL から転送先の内容を取得する場合に必要となる。
\item[\code{-f}] HTTP 400 以上の応答を受けた場合に、終了状態を失敗とする。スクリプトから終了状態を検査すれば、エラーページを正常な取得結果として後続処理へ渡すことを防げる。
\item[\code{-sS}] 進捗表示を抑制し、エラーメッセージは表示する。スクリプト内で出力を制御する場合に用いる。
\end{description}

\subsection{\code{wget}}

\code{wget}\index{wget@\code{wget}|textbf} は、URL を指定すると既定でファイルへ保存する。\code{-c} を指定すれば、サーバーが範囲指定による取得に対応している場合に、途中まで保存したファイルの続きから取得できる。

\begin{lstlisting}[numbers=none, xleftmargin=2\zw]
$ wget https://example.com/data/run001.csv
$ wget -O events.csv https://example.com/data/run001.csv
$ wget -c https://example.com/archive.tar.gz
\end{lstlisting}

1 行目では、URL の末尾にある \code{run001.csv} という名前で保存される。2 行目の \code{-O} は保存名を指定する。3 行目の \code{-c} は、途中まで保存されたファイルがある場合に、その続きから取得を再開する。

\begin{description}
\item[\code{-O}] 保存名を自分で指定する。既定の名前を使いたくないときに使う。
\item[\code{-c}] 途中まで保存したファイルの続きから取得を試みる。サーバーが範囲指定による取得に対応していなければ、先頭から取得し直す場合がある。
\item[\code{-q}] 通常のメッセージを表示しない。
\item[\code{--show-progress}] quiet モードを指定していても、進捗表示を有効にする。
\end{description}

\code{wget} が手元にない環境でも、\code{curl -O} または \code{curl -o} でファイルへ保存できる。受信内容をパイプで次のコマンドへ渡す場合は、既定で標準出力へ送る \code{curl} をオプションなしで利用できる。

\section{gzip・tar・zip による圧縮と展開}

\subsection{\code{gzip}: 1 ファイルの圧縮}

\code{gzip}\index{gzip@\code{gzip}|textbf} は、1 つのファイルを gzip 形式の 1 つの圧縮ファイルへ変換する。ディレクトリ構造や複数のファイルを一つにまとめる機能は持たない。

\begin{lstlisting}[numbers=none, xleftmargin=2\zw]
$ gzip -k events.csv
$ ls -lh events.csv*
$ gzip -cd events.csv.gz | head -n 2
$ gzip -cd events.csv.gz > events_restored.csv
\end{lstlisting}

1 行目の \code{-k} は元ファイルを残す指定である。付けない場合は、通常 \code{events.csv} が \code{events.csv.gz} に置き換わる。2 行目では圧縮前後のサイズを見比べている。3 行目の \code{-cd} は、展開した内容をファイルへ戻さず標準出力へ流す指定であり、\code{head} や \code{grep} と組み合わせられる。4 行目は同じ出力をリダイレクトし、元の \code{events.csv} を残したまま \code{events\_restored.csv} へ展開する。\code{gunzip events.csv.gz} で元の名前へ戻すには、同名の \code{events.csv} が存在しないことが必要である。

\begin{description}
\item[\code{-k}] 圧縮または展開の完了後も入力ファイルを残す。
\item[\code{-d}] gzip 形式のファイルを展開する。
\item[\code{-c}] 圧縮または展開した結果を標準出力へ送る。ファイルを作らずに \code{head} や \code{grep} へ渡せる。
\item[\code{-1} から \code{-9}] 圧縮にかける時間と出力サイズの兼ね合いを指定する。数字を大きくすると、通常は処理時間が長くなる代わりに出力ファイルが小さくなる。実際の差は入力内容に依存するため、必要な場合に同じ入力で処理時間と容量を比較する。
\end{description}

\code{gzip} 自体は複数ファイルを 1 つにまとめない。複数ファイルやディレクトリ構造を保持する場合は、次の \code{tar} でアーカイブを作成し、そのアーカイブを gzip で圧縮する。

\subsection{\code{tar}: tar アーカイブの作成と gzip 圧縮}

\code{tar}\index{tar@\code{tar}|textbf} は、複数のファイルやディレクトリ構造を一つのファイルへ格納する。この格納先を tar アーカイブという。\code{-z} を指定すると、作成した tar アーカイブを gzip で圧縮した \code{.tar.gz} ファイルを作る。ここで用いる \code{-c}、\code{-t}、\code{-x}、\code{-f}、\code{-z} の意味は GNU tar 公式マニュアル~\cite{gnu_tar_manual} に基づく。

\begin{lstlisting}[numbers=none, xleftmargin=2\zw]
$ tar -czf archive.tar.gz raw/ notes.txt
$ tar -tzf archive.tar.gz
$ tar -xzf archive.tar.gz
\end{lstlisting}

1 行目は、\code{raw/} と \code{notes.txt} を \code{archive.tar.gz} へ格納し、gzip で圧縮する。2 行目の \code{-t} は、展開せずに収録項目を表示する。3 行目の \code{-x} は内容を展開する。
\begin{description}
\item[\code{-c}] 新しくアーカイブを作る。
\item[\code{-t}] 中身一覧だけを見る。
\item[\code{-x}] 展開する。
\item[\code{-z}] \code{gzip} 圧縮を併用する。
\item[\code{-f}] 直後の文字列をアーカイブ名として扱う。
\end{description}

\code{tar -czf archive.tar.gz ...} の \code{-c} は tar アーカイブの作成、\code{-z} は gzip 圧縮、\code{-f} は直後の \code{archive.tar.gz} をアーカイブ名として扱う指定である。内容一覧には \code{-t}、展開には \code{-x} を用いる。

\subsection{\code{zip}: ZIP アーカイブの作成と展開}

\code{zip}\index{zip@\code{zip}|textbf} は、複数のファイルやディレクトリを圧縮し、1 つの ZIP アーカイブへ格納する。\code{tar} とはオプションと展開コマンドが異なるため、ファイル名の拡張子に対応するコマンドを選ぶ。

\begin{lstlisting}[numbers=none, xleftmargin=2\zw]
$ zip -r archive.zip raw/ notes.txt
$ unzip -l archive.zip
$ unzip archive.zip
\end{lstlisting}

1 行目の \code{-r} は、指定したディレクトリ以下を再帰的に格納する。2 行目の \code{unzip -l} は展開せずに収録項目を表示し、3 行目は内容を展開する。形式に応じて、\code{.zip} には \code{unzip}、\code{.tar.gz} には \code{tar -xzf}、単一の \code{.gz} ファイルには \code{gunzip} または \code{gzip -d} を用いる。

\section{環境変数と PATH}

\subsection{環境変数の定義と参照}

シェルやプログラムには、「名前」と「値」の組として渡される設定がある。これが環境変数\index{かんきょうへんすう@環境変数|textbf}である。環境変数は、シェルの中だけで完結するものではなく、そこから起動した Python や各種コマンドにも引き継がれる。現在設定されている主な環境変数の値は、次のコマンドで表示できる。

\begin{lstlisting}[numbers=none, xleftmargin=2\zw]
$ echo $HOME
$ echo $SHELL
$ echo $PATH
\end{lstlisting}

実行結果は、たとえば次のようになる。

\begin{lstlisting}[numbers=none, xleftmargin=2\zw]
/Users/student
/bin/zsh
/opt/homebrew/bin:/usr/local/bin:/usr/bin:/bin:/usr/sbin:/sbin
\end{lstlisting}

1 行目はホームディレクトリ、2 行目はログインシェル、3 行目は実行ファイルを探索するディレクトリの一覧である。\code{PATH}\index{path@\code{PATH}|textbf} の値はコロン \code{:} で区切られ、シェルは左から順に同名の実行ファイルを探索する。\code{PYENV\_ROOT} や \code{DATA\_ROOT} も環境変数であり、起動するプログラムへ名前付きの値として引き継がれる。新たな環境変数は次の形で設定する。

\begin{lstlisting}[numbers=none, xleftmargin=2\zw]
$ export DATA_ROOT="$HOME/data/events"
$ echo $DATA_ROOT
\end{lstlisting}

結果の例は次の通りである。

\begin{lstlisting}[numbers=none, xleftmargin=2\zw]
/Users/student/data/events
\end{lstlisting}

\code{export} 自体は成功しても通常は何も表示しない。続けて \code{echo \$DATA\_ROOT} を実行し、意図した値が入ったことを確認する。ここでは \code{\$HOME} が先に展開され、その結果として絶対パスが保存されている。

ただし、この設定は通常そのターミナルを閉じると消える。毎回使う値なら、対話シェルが読み込む設定ファイルに \code{export ...} を書いておく。第~\ref{SEC:shell}~章で述べたように、\code{echo \$SHELL} でログインシェルを、\code{ps -p \$\$ -o comm=} で現在のシェルを確認できる。zsh では通常 \code{\textasciitilde/.zshrc}、bash では環境に応じて \code{\textasciitilde/.bashrc} または \code{\textasciitilde/.bash\_profile} に追記する。

\subsection{展開と引用}

シェルはコマンドの実行前に、\code{\$HOME} のような変数展開と、\code{\$(pwd)} のようなコマンド置換を行う。展開の有無と、展開後の値を一つの引数として保つかどうかは引用符によって異なる。

\begin{lstlisting}[numbers=none, xleftmargin=2\zw]
$ echo $HOME
$ echo "$HOME"
$ echo '$HOME'
$ name="analysis result"
$ echo $name
$ echo "$name"
$ echo "$(pwd)"
$ echo "`pwd`"
\end{lstlisting}

1 行目と 2 行目は、どちらも \code{\$HOME} を展開する。\code{HOME} に空白がなければ表示は同じになる。しかし、クォートなしの変数展開は、展開後にフィールド分割とファイル名展開の対象になる。\code{"\$HOME"} は展開結果を一つの引数として保つため、文字列として渡す変数は原則としてダブルクォートで囲む。

3 行目のシングルクォート \code{'...'} 内では変数展開が行われないため、\code{\$HOME} という文字列がそのまま表示される。\code{name="analysis result"} の後に引用符なしで \code{echo \$name} と書くと、展開後の空白によって 2 引数へ分割される。\code{echo "\$name"} では空白を含む値が一つの引数として渡る。\code{\$(pwd)} とバッククォートは、いずれも \code{pwd} の標準出力によるコマンド置換を表す。\code{\$(...)} は入れ子を明示できるため、本書ではこちらを用いる。

この違いは、設定ファイルへ 1 行追記するときに特に重要である。たとえば \code{echo 'export PATH="\$HOME/bin:\$PATH"' >> \textasciitilde/.zshrc} のようにシングルクォートで囲むと、\code{\$HOME} や \code{\$PATH} を展開せず、その文字列をそのまま \code{\textasciitilde/.zshrc} へ追記できる。もしダブルクォートで囲むと、その場で現在の値へ展開されて書き込まれるので、意図とずれることがある。

また、空白と改行の扱いにも注意が必要である。クォートの外では、空白やタブや改行が単語の区切りになる。そのため、空白を含むファイル名はクォートしないと壊れる。

\begin{lstlisting}[numbers=none, xleftmargin=2\zw]
$ mkdir "analysis result"
$ cp "analysis result/data.txt" backup/
\end{lstlisting}

1 つの長いコマンドを複数行へ分けたいときは、行末に \code{\textbackslash} を置いて改行を継続する。

\begin{lstlisting}[numbers=none, xleftmargin=2\zw]
$ python analyze.py \
  --input data/events.txt \
  --output result.txt
\end{lstlisting}

引用符や括弧の外にある改行は、通常はコマンドの終端として解釈される。行末の \code{\textbackslash} は直後の改行を除去し、次の物理行を同じコマンドとして継続する。

\subsection{PATH と実行ファイル}

シェルが \code{python3} や \code{ls} を実行できるのは、それらが実行ファイルとして保存されており、しかもシェルがその場所を知っているからである。その探索に使う環境変数が \code{\$PATH} である。実際に実行ファイルがどこから使われているかを探すには、次のコマンドを用いる。

\begin{lstlisting}[numbers=none, xleftmargin=2\zw]
$ command -v python3
$ command -v ls
$ echo $PATH
\end{lstlisting}

出力の例は次のようになる。

\begin{lstlisting}[numbers=none, xleftmargin=2\zw]
/opt/homebrew/bin/python3
/bin/ls
/opt/homebrew/bin:/usr/local/bin:/usr/bin:/bin:/usr/sbin:/sbin
\end{lstlisting}

\code{command -v python3} は、現在のシェルが \code{python3} という名前をどのコマンドに解決するかを表示する。例のようにパスが表示される場合もあれば、エイリアスやシェル関数として定義されている場合もある。\code{PATH} で \code{/opt/homebrew/bin} が \code{/usr/bin} より左にあるため、同名の実行ファイルが両方にあれば前者が優先される。Python 自身が認識する実行ファイルは、\code{python3 -c "import sys; print(sys.executable)"} で別に確認できる。

\code{bin} は実行ファイルを置くディレクトリに用いられる慣習的な名前である。仮想環境の \code{.venv/bin} にも Python と、その環境へインストールしたコマンドの起動ファイルが置かれる。

独自の実行ファイルを置く \code{\textasciitilde/my\_scripts} をコマンド探索の対象にする例を示す。既存の探索先を保持するために現在の \code{\$PATH} を右辺へ含め、新しいディレクトリを先頭へ追加する。対話シェルの起動時に同じ設定を適用する場合は、利用中のシェルが読み込む設定ファイルに記述する。

\begin{lstlisting}[numbers=none, xleftmargin=2\zw]
$ export PATH="$HOME/my_scripts:$PATH"
\end{lstlisting}

\section{リダイレクトとログ保存}

\subsection{\code{echo}、標準出力、標準エラー}

シェルのコマンドは、単に画面へ文字を出すだけではなく、出力先を切り替えられる。ここで区別したいのは、通常の結果を流す標準出力（standard output, stdout）と、エラーメッセージを流す標準エラー（standard error, stderr）である。

\code{echo} は、引数として受け取った文字列を標準出力へ送る。

\begin{lstlisting}[numbers=none, xleftmargin=2\zw]
$ echo "hello"
$ echo "$HOME"
\end{lstlisting}

1 行目は \code{hello}、2 行目は変数展開後の \code{\$HOME} の値を標準出力へ送る。端末上では標準出力が画面へ接続されているが、次節のリダイレクトによってファイルへ変更できる。

\subsection{\code{>} と \code{>>}}

\code{>} は標準出力をファイルへ書き出す。既存のファイルがあれば中身を上書きする。

\begin{lstlisting}[numbers=none, xleftmargin=2\zw]
$ echo "alpha" > notes.txt
$ cat notes.txt
alpha
\end{lstlisting}

この例では、\code{echo} の結果が画面ではなく \code{notes.txt} に保存される。\code{notes.txt} がすでに存在する場合、その内容は新しい出力で置き換えられる。既存ファイルを保持する必要がある場合は、出力先を確認してから実行する。

\code{>>} は追記であり、既存の内容を残したまま末尾に足す。

\begin{lstlisting}[numbers=none, xleftmargin=2\zw]
$ echo "alpha" > notes.txt
$ echo "beta" >> notes.txt
$ cat notes.txt
alpha
beta
\end{lstlisting}

\code{>>} は既存ファイルを保持して末尾へ追記する。\code{>} は既存内容を置き換えるため、設定ファイルへ行を追加する操作には \code{>>} を用いる。いずれの場合も、実行前に対象パスを確認する。

\subsection{\code{1>} と \code{2>}}

\code{>} は実際には \code{1>} と同じ意味であり、標準出力だけをファイルへ送る。標準エラーを別ファイルへ分けたいときは \code{2>} を使う。

\begin{lstlisting}[numbers=none, xleftmargin=2\zw]
$ python analyze.py > output.log
$ python analyze.py 1> output.log 2> error.log
\end{lstlisting}

1 行目では通常の出力だけが \code{output.log} に入り、エラーメッセージは画面に残る。2 行目では、通常の出力を \code{output.log} に、エラーメッセージを \code{error.log} に分けて保存している。長時間の処理では、両者を分けると失敗原因を追跡しやすい。

標準出力と標準エラーを同じファイルへまとめたいときは、\code{2>\&1} を使う。

\begin{lstlisting}[numbers=none, xleftmargin=2\zw]
$ python analyze.py > run.log 2>&1
\end{lstlisting}

これは、標準エラーの行き先を標準出力と同じ先へ合わせる、という意味である。ログを 1 本にまとめたいときに使う。

\subsection{\code{tee} による画面表示と保存}

\code{>} や \code{>>} で出力をファイルへ送ると、その出力は画面に表示されない。\code{tee}\index{tee@\code{tee}|textbf} は標準入力を標準出力へ複製すると同時に、指定したファイルへ保存する。

\begin{lstlisting}[numbers=none, xleftmargin=2\zw]
$ python analyze.py | tee stdout.log
$ python analyze.py 2>&1 | tee combined.log
\end{lstlisting}

1 行目は標準出力だけを画面と \code{stdout.log} の両方へ流す。2 行目は標準エラーもまとめ、画面と \code{combined.log} へ同じ内容を出力する。二つの行は保存対象の違いを示す代替例である。

\code{tee} に \code{-a} を付けると追記になる。

\begin{lstlisting}[numbers=none, xleftmargin=2\zw]
$ python analyze.py 2>&1 | tee -a run.log
\end{lstlisting}

この場合は、既存の \code{run.log} を消さずに末尾へ足していく。毎日の実行結果を 1 つのログへ蓄積したい場合に向く。

\section{権限とスクリプトの実行}

\subsection{ファイル権限の表記と解釈}
\label{subsec:mod}
ファイルにはアクセス権\index{あくせすけん@アクセス権|textbf}がある。読み取れるか、書き込めるか、実行できるかを \code{rwx} で表す。\code{ls -l} を使うと確認しやすい。

\begin{lstlisting}[numbers=none, xleftmargin=2\zw]
-rwxr-xr-x  1 user  staff  1234 Apr  8 10:00 run_analysis.sh
-rw-r--r--  1 user  staff   411 Apr  8 10:05 notes.txt
-rw-------  1 user  staff   411 Apr  8 10:05 id_ed25519
\end{lstlisting}

最初の 10 文字のうち、先頭 1 文字はファイル種別であり、残り 9 文字が権限である。3 文字ずつに区切って、所有者、グループ、その他の利用者に対する権限を表す。
文字と権限の対応は表~\ref{table:mod}のとおりである。

権限の変更には、\code{chmod}\index{chmod@\code{chmod}|textbf} コマンドを用いる。
所有者（\code{user}）に実行権限（\code{x}）を加える指定は \code{u+x} である。その他の利用者（\code{other}）から読み取り権限（\code{r}）と書き込み権限（\code{w}）を除く指定には \code{o-rw} を用いる。グループ（\code{group}）とその他の利用者の権限を読み取りだけにする指定は \code{go=r} である。\code{a+x} の \code{a}（\code{all}）は、すべての利用者を対象とする。

\begin{lstlisting}[numbers=none, xleftmargin=2\zw]
$ chmod u+x my_script.sh
$ chmod o-rw notes.txt
$ chmod go=r read_only.txt
$ chmod a+x run_analysis.sh
\end{lstlisting}

\begin{table}[hbtp]
\centering
\caption{権限の記号と数字}
\label{table:mod}
\begin{tabular}{lll}
\toprule
記号 & 数字 & 意味 \\
\midrule
\code{r} & 4 & 読み取り \\
\code{w} & 2 & 書き込み \\
\code{x} & 1 & 実行 \\
\bottomrule
\end{tabular}
\end{table}

表~\ref{table:mod}に示すように、権限は数字でも指定できる。各桁は読み取りを 4、書き込みを 2、実行を 1 として合計した値である。たとえば \code{7} は \code{rwx}、\code{6} は \code{rw-}、\code{5} は \code{r-x} に対応する。したがって、\code{755} は所有者が \code{rwx}、グループとその他の利用者が \code{r-x}、\code{644} は所有者が \code{rw-}、グループとその他の利用者が \code{r--} である。SSH の秘密鍵には、所有者だけに読み書きを許可する \code{600} を指定する。設定例を次に示す。

\begin{lstlisting}[numbers=none, xleftmargin=2\zw]
$ chmod 755 run_analysis.sh
$ chmod 644 notes.txt
$ chmod 600 ~/.ssh/id_ed25519
\end{lstlisting}

\subsection{スクリプトの直接実行}

Python スクリプトは \code{python3 hello.py} の形で実行するなら、実行権限がなくてもよい。これに対して \code{./hello.py} のように直接実行する場合には、どのプログラムで解釈するかをファイル先頭に書いておく必要がある。
その先頭行の \code{\#!} で始まる記法を、シバン (shebang) という。

\begin{lstlisting}[language=Python]
#!/usr/bin/env python3

print("hello")
\end{lstlisting}

ファイルを直接実行するには、シバンに加えて実行権限も必要である。次の \code{u+x} は、所有者（user）に実行権限を加える指定である。

\begin{lstlisting}[numbers=none, xleftmargin=2\zw]
$ chmod u+x hello.py
$ ./hello.py
\end{lstlisting}

\code{\#!/usr/bin/env python3} は、\code{env} が \code{PATH} から \code{python3} を探し、その実行ファイルでスクリプトを解釈する指定である。ファイルを \code{./script.py} として直接実行するには、適切なシバンと実行権限の両方が必要になる。付録~\ref{app:shell_scripting} では、同じ仕組みに基づくシェルスクリプトと、\code{for} ループや条件分岐による一括処理を扱う。

\section{USGS カタログの取得・調査・圧縮}

各課題の記録対象は、コマンドを実行した事実だけでなく、入力、主要なオプション、出力の読み方までを含む。ネットワーク上のデータを取得できない場合でも、1 問目と、手元の任意の CSV を使った 3 問目は実行できる。

\begin{enumerate}
\item \textbf{章末問題用の作業ディレクトリの作成}。
次のコマンドを実行し、どのようなディレクトリが作成されるか確認せよ。
\begin{lstlisting}[numbers=none, xleftmargin=2\zw]
$ mkdir -p earthquake_practice/{raw,compressed,scripts,figures,results}
$ cd earthquake_practice
$ pwd
$ ls
\end{lstlisting}
シェルは \code{earthquake\_practice/\{raw,compressed,scripts,figures,results\}} を、コンマで区切られた 5 通りのパスへ展開してから \code{mkdir} へ渡す。この機能をブレース展開という。\code{-p} は親ディレクトリも作り、対象が既存でもエラーにしない。\code{ls} の結果に 5 個の名前が現れることを確認し、\code{pwd} の出力がどのディレクトリを示すかも説明せよ。

\item \textbf{USGS 地震カタログの取得とシェルによる内容確認}。
本書の章末課題では、アメリカ地質調査所（\textbf{U}nited \textbf{S}tates \textbf{G}eological \textbf{S}urvey、\textbf{USGS}）が公開している地震カタログを使用する。
以下のコマンドを実行し、2024年1月のカタログを \code{raw/} へ保存し、行数、先頭数行、末尾数行、特定の地名を含む行を調べよ。
\begin{lstlisting}[numbers=none, xleftmargin=2\zw]
$ curl -L "https://earthquake.usgs.gov/fdsnws/event/1/query?format=csv&starttime=2024-01-01&endtime=2024-01-31T23:59:59.999999&limit=20000&orderby=time-asc" \
  -o raw/usgs_global_2024_01.csv
$ ls -lh raw/usgs_global_2024_01.csv
$ wc -l raw/usgs_global_2024_01.csv
$ head -n 3 raw/usgs_global_2024_01.csv
$ tail -n 3 raw/usgs_global_2024_01.csv
$ grep -i "Japan" raw/usgs_global_2024_01.csv | head
\end{lstlisting}
どのコマンドが「件数確認」「冒頭確認」「末尾確認」「文字列検索」に対応しているかを整理してまとめよ。
\code{curl} の \code{-L} は転送先が変更された場合に追従し、\code{-o} は保存先を指定する。\code{wc -l} の値にはヘッダ行も含まれるので、データ件数との違いを説明せよ。\code{grep -i} の \code{-i} と、後段の \code{head} が検索結果を何件に制限するかも記録せよ。

\item \textbf{同じ CSV に対する 3 種類の圧縮形式の比較}。
\code{gzip}、\code{tar.gz}、\code{zip} をそれぞれ 1 つずつ作り、圧縮前後の容量を確認せよ。
\begin{lstlisting}[numbers=none, xleftmargin=2\zw]
$ gzip -c raw/usgs_global_2024_01.csv > compressed/usgs_global_2024_01.csv.gz
$ tar -czf compressed/usgs_global_2024_01.tar.gz raw/usgs_global_2024_01.csv
$ zip compressed/usgs_global_2024_01.zip raw/usgs_global_2024_01.csv
$ gzip -cd compressed/usgs_global_2024_01.csv.gz | head -n 2
$ tar -tzf compressed/usgs_global_2024_01.tar.gz
$ unzip -l compressed/usgs_global_2024_01.zip
$ wc -c raw/usgs_global_2024_01.csv compressed/*
\end{lstlisting}
三つの形式に同じ CSV が入っていることは、\code{gzip -cd}、\code{tar -tzf}、\code{unzip -l} の出力で確認する。\code{gzip -c} は元の CSV を変更せず、圧縮結果を標準出力へ書き出す。\code{tar -czf} は指定した CSV を tar アーカイブに格納し、そのアーカイブを gzip で圧縮する。\code{zip} も同じ CSV を ZIP アーカイブに格納する。今回の入力は単一ファイルであるため、ディレクトリを再帰的にたどる \code{zip -r} は使わない。

\code{wc -c} は各ファイルのバイト数を表示する。その値から、\code{容量比 = 圧縮後のバイト数 / 元 CSV のバイト数} をそれぞれ求める。この値が小さいほど、圧縮後の容量が元の容量より小さい。コマンドの各オプションの役割も、本節の説明と対応付けて記録する。

\item \textbf{発展課題}。
URL の構造について調べよ。
カタログの URL（\url{https://earthquake.usgs.gov/fdsnws/event/1/query?format=csv&starttime=2024-01-01&endtime=2024-01-31T23:59:59.999999&limit=20000&orderby=time-asc}）を確認し、\code{https}、\code{earthquake.usgs.gov}、\code{usgs.gov}、\code{.gov}、\code{/fdsnws/event/1/query}、\code{?format=csv\&starttime=2024-01-01\&endtime=2024-01-31T23:59:59.999999\&limit=20000\&orderby=time-asc} のそれぞれについて、名称や役割をまとめよ。
ここでは順に、スキーム、ホスト名、登録ドメイン、トップレベルドメイン、パス、クエリ文字列を表す。クエリ文字列は \code{\&} で複数の \code{名前=値} を区切る。\code{format}、\code{starttime}、\code{endtime}、\code{limit}、\code{orderby} が取得結果をどう変えるかを、2 問目の出力と対応させて説明せよ。
\end{enumerate}

シェルでファイルの所在、内容、権限、入出力を確認できれば、Python の導入後も実行されるプログラムと保存先を追跡できる。次章では、\code{PATH} と実行ファイルの知識に基づく Python のバージョン、仮想環境、プロジェクト依存関係の管理を扱う。
